\chapter{المحددات والجمل الخطية}\label{ch:b1:det}

يضغط المحدد في سلّم واحد الجوابَ عن سؤال «هل هذه
المتجهات $n$ \omterm{def:b1:vspaces:free}{أساس}؟» --- ويقيس، هندسيًا، الحجم
الذي تولّده. ونميّزه بخصائصه (متعدد الخطية، متناوب،
منظَّم)، ونحسبه في البُعدين $2$ و $3$ وبالنشر
بالعوامل المرافقة عمومًا، ونشغّله على الجمل الخطية،
إلى جانب الخوارزمية الصالحة لكل غرض: \omterm{met:b1:det:gauss}{إزاحة غاوس}.

\section{المحدد}

\begin{theorem}[التمييز]\label{thm:b1:det:def}
يوجد \omterm{def:b1:logic:map}{تطبيق} واحد بالضبط $\det \colon \mathcal{M}_n(K) \to K$،
منظورًا إليه دالةً في الأعمدة $n$، يكون:
\begin{enumerate}
  \item \emph{\omterm{def:b1:linmaps:def}{خطيًا} في كل عمود} (والبقية مثبَّتة)؛
  \item \emph{متناوبًا}: أي إن مبادلة عمودين تغيّر الإشارة
        (ومنه فالعمودان المتساويان يعطيان $0$)؛
  \item \emph{منظَّمًا}: $\det I_n = 1$.
\end{enumerate}
ومن أجل $n = 2$ و $3$:\index{محدد}
\[
\begin{vmatrix} a & b\\ c & d\end{vmatrix} = ad - bc,
\qquad
\begin{vmatrix} a & b & c\\ d & e & f\\ g & h & i\end{vmatrix}
= aei + bfg + cdh - ceg - bdi - afh
\]
(وقاعدة ساروس للمقدار $3 \times 3$: جداءات الأقطار النازلة
ناقص الصاعدة).
\end{theorem}
\admitted

\begin{remark}
من أجل $n = 2$: يعطي النشر بثنائية الخطية على الأعمدة القانونية
الصيغةَ، وهي بالعكس تحقق البديهيات --- أي برهان كامل؛ و $n = 3$
مطابق مع حدود أكثر. وأمّا الحالة العامة
(الوجود عبر المجموع على التبديلات، والوحدانية بالنشر نفسه)
فتقتضي توقيع \omterm{def:b1:counting:objects}{تبديلة} وهي مؤجَّلة إلى السنة الثانية؛ ونستعمل
بحرية البديهيات والنتائج
أدناه.

وإليك نشر $n = 2$ كاملًا، لأنه القالب: فمع
العمودين $C_1 = a\,e_1 + c\,e_2$ و $C_2 = b\,e_1 + d\,e_2$،
تعطي ثنائية الخطية
\[
\det(C_1, C_2) = ab\det(e_1, e_1) + ad\det(e_1, e_2)
+ cb\det(e_2, e_1) + cd\det(e_2, e_2),
\]
ويقتل التناوب الأزواج المكررة بينما يقلب
$\det(e_2, e_1) = -\det(e_1, e_2)$: فينهار \omterm{def:b1:logic:map}{التطبيق} كله إلى
$(ad - bc)\det(e_1, e_2) = ad - bc$ بالتنظيم. والوحدانية
مرئية في الحساب نفسه --- إذ لم تترك البديهيات أيّ
اختيار عند أيّ خطوة --- وهذه بالضبط واقعة الوحدانية المسلَّمة
المستعملة في برهان قاعدة الجداء أدناه.
\end{remark}

\begin{theorem}[الخصائص]\label{thm:b1:det:props}
من أجل $A, B \in \mathcal{M}_n(K)$:
\begin{enumerate}
  \item إضافة مضاعف لعمود إلى عمود آخر لا تغيّر
        المحدد؛ وضرب عمود في $\lambda$ يضربه
        في $\lambda$ (ومنه $\det(\lambda A) = \lambda^n \det A$)؛
  \item $\det(AB) = \det A\, \det B$؛
  \item تكون $A$ قابلة للقلب $\iff$ $\det A \neq 0$ $\iff$ تكوّن الأعمدة
        \omterm{def:b1:vspaces:free}{أساسًا} للمقدار $K^n$؛ وعندئذ $\det(A^{-1}) = (\det
        A)^{-1}$؛
  \item $\det(A^{\mathsf T}) = \det A$ --- ومنه فكل قاعدة على الأعمدة
        قاعدةٌ على السطور كذلك؛
  \item محدد مصفوفة مثلثية جداءُ مركّباتها
        القطرية.
\end{enumerate}
\end{theorem}

\begin{proof}
(1) بالخطية، $\det(\dots, C_i + \lambda C_j, \dots) = \det A +
\lambda\det(\dots, C_j, \dots)$ حيث للمحدد الثاني عمودان
متساويان: فهو معدوم.

(2) ثبّت $A$ وانظر في $\varphi(B) = \det(AB)$ دالةً في
أعمدة $B$: ولأن للمقدار $AB$ الأعمدة $AB_j$، تكون $\varphi$
متعددة الخطية ومتناوبة في المقادير $B_j$. ونقبل، مع
\cref{thm:b1:det:def}، \omterm{def:b1:logic:statement}{عبارة} وحدانيته في صورتها المسلَّمة:
\emph{كل} \omterm{def:b1:logic:map}{تطبيق} متعدد الخطية متناوب $\varphi$ للأعمدة
يساوي $\varphi(I_n) \cdot \det$. وهنا $\varphi(I_n) = \det A$، ومنه
$\det(AB) = \det A \cdot \det B$.

(3) إذا كانت $A$ قابلة للقلب: $\det A\,\det A^{-1} = \det I = 1$، ومنه
$\det A \neq 0$ وتصحّ صيغة المعكوس. وإذا لم تكن $A$
قابلة للقلب، فأعمدتها مرتبطة (\cref{cor:b1:linmaps:samedim}
و \cref{prop:b1:linmaps:basis})؛ وبالتعبير عن عمود بالبقية
والنشر بالخطية تبقى محددات ذات عمودين
متساويين: ومنه $\det A = 0$. \omterm{def:b1:logic:statement}{وعبارة} \omterm{def:b1:vspaces:free}{الأساس} هي
\cref{prop:b1:findim:twoofthree}.

(4) مقبولة مع البناء العام (فهي مباشرة على
صيغة التبديلات)؛ ونسجّلها لاستعمال \omterm{met:b1:matrices:gauss}{العمليات على السطور}.

(5) إذا انعدمت مركّبة قطرية ما، فإن الأعمدة $k$ الأولى
مرتبطة من أجل $k$ ما (باعتبارات الرتبة) و $\det = 0 =$ الجداء.
وإلا فنظّف كل عمود تحت وإلى اليسار بعمليات من النوع
(1) --- وهذا ممكن في الشكل المثلثي --- فتُبلغ المصفوفة
القطرية، ومحددها جداء المركّبات
بتعدد الخطية انطلاقًا من $I_n$.
\end{proof}

\begin{example}[القواعد، متحققًا منها على أعداد]
\label{ex:b1:det:rulescheck}
خذ $A = \begin{pmatrix} 1 & 2\\ 3 & 4\end{pmatrix}$ ($\det A =
-2$) و $B = \begin{pmatrix} 0 & 1\\ 1 & 1\end{pmatrix}$ ($\det
B = -1$). عندئذ
\[
AB = \begin{pmatrix} 2 & 3\\ 4 & 7\end{pmatrix},
\quad \det(AB) = 14 - 12 = 2 = (-2)(-1) ;
\qquad
\det(A^{\mathsf T}) = \begin{vmatrix} 1 & 3\\ 2 & 4
\end{vmatrix} = -2 = \det A .
\]
فتتأكد الضربية والحفظ بالنقل --- بينما
تفشل الجمعية \emph{الخاطئة} على الزوج نفسه:
\[
\det(A + B) = \begin{vmatrix} 1 & 3\\ 4 & 5\end{vmatrix} = -7
\neq \det A + \det B = -3 .
\]
وثلاثون ثانية من حساب من هذا النوع، بعد استدعاء أيّ متطابقة
محددات، هي أرخص تأمين متوفر ضد الأخطاء.
\end{example}

\begin{example}[المحددات مساحاتٍ]
\label{ex:b1:det:area}
متوازي الأضلاع المولَّد بالمقدارين $u = (2, 0)$ و $v = (1, 3)$ له
قاعدة $2$ وارتفاع $3$: فمساحته $6$. و
\[
\begin{vmatrix} 2 & 1\\ 0 & 3\end{vmatrix} = 6 :
\]
فالمحدد $2\times2$ \emph{هو} المساحة الموقَّعة
لمتوازي أضلاع عموديه. وتعيد البديهيات سرد الهندسة:
فإضافة مضاعف لعمود إلى الآخر \emph{قصّ}،
ينزلق بمتوازي الأضلاع موازيًا لضلع دون تغيير
القاعدة أو الارتفاع (العملية~(1) في \cref{thm:b1:det:props})؛
وضرب عمود في سلّم يضرب المساحة؛ ومبادلة العمودين تقلب
التوجيه، ومنه الإشارة، $\det(v, u) = -6$. وفي $\R^3$
تعطي القراءة نفسها حجومًا موقَّعة، ويصير $\abs{\det}$
عاملَ تسليم الحجم الشامل للتطبيقات الخطية --- وهي الواقعة
وراء صيغة تغيير المتغيرات في التكاملات المضاعفة
في مجلّد السنة~الثانية.
\end{example}

\begin{theorem}[النشر بالعوامل المرافقة]\label{thm:b1:det:cofactor}
لتكن $A \in \mathcal{M}_n(K)$ وليكن $\Delta_{ij}$ محدد
$A$ بعد حذف السطر $i$ والعمود $j$. عندئذ، من أجل أيّ عمود
مثبَّت $j$ (أو سطر، بالنقل):\index{نشر بالعوامل المرافقة}
\[
\det A = \sum_{i=1}^{n} (-1)^{i+j}\, a_{ij}\, \Delta_{ij} .
\]
\end{theorem}
\admitted

\begin{example}\label{ex:b1:det:cofactor}
بالنشر على امتداد العمود الأول:
\[
\begin{vmatrix}
2 & 1 & 0\\
1 & 2 & 1\\
0 & 1 & 2
\end{vmatrix}
= 2\begin{vmatrix} 2 & 1\\ 1 & 2\end{vmatrix}
- 1\begin{vmatrix} 1 & 0\\ 1 & 2\end{vmatrix}
= 2 \times 3 - 2 = 4 .
\]
والاستراتيجية: أنشئ أصفارًا أولًا (\omterm{met:b1:matrices:gauss}{بعمليات على السطور} أو الأعمدة)، ثم انشر
على امتداد أخلى سطر أو عمود.
\end{example}

\begin{example}[معكوس العوامل المرافقة، مرة واحدة باليد]
\label{ex:b1:det:cofactorinverse}
من أجل $A = \begin{pmatrix} 1 & 1 & 0\\ 0 & 1 & 1\\ 1 & 0 &
1\end{pmatrix}$: $\det A = 1(1) - 1(-1) + 0 = 2$. وتتجمّع
العوامل المرافقة التسعة $(-1)^{i+j}\Delta_{ij}$ في
\[
\operatorname{Com}(A) = \begin{pmatrix}
1 & 1 & -1\\
-1 & 1 & 1\\
1 & -1 & 1
\end{pmatrix},
\qquad
A^{-1} = \frac{1}{\det A}\operatorname{Com}(A)^{\mathsf T}
= \frac12\begin{pmatrix}
1 & -1 & 1\\
1 & 1 & -1\\
-1 & 1 & 1
\end{pmatrix},
\]
وهي الصيغة المذكورة في \cref{exo:b1:det:8}. وتحقق من زوج
سطر-عمود واحد: (السطر $1$ من $A$)(العمود $1$ من $A^{-1}$) $= \frac12(1 + 1
+ 0) = 1$، وإزاء العمود $2$: $\frac12(-1 + 1 + 0) = 0$.
أي تسعة محددات $2\times2$ من أجل معكوس $3\times3$ واحد: فأصلًا
عند هذا الحجم، يكون تقليص السطور (\cref{exo:b1:det:3}) أرخص ---
وقيمة صيغة العوامل المرافقة نظرية (الصحيحية في
\cref{exo:b1:det:8}، وقابلية اشتقاق المعكوس في مجلّدات لاحقة)،
لا حسابية.
\end{example}

\begin{example}[قاعدة الكتل المثلثية، في الحجم $4$]
\label{ex:b1:det:blocktriangular}
الادعاء: $\det\begin{pmatrix} M & N\\ 0 & P\end{pmatrix} = \det
M\,\det P$ من أجل كتل $2\times2$. نظّف كتلة $N$ بعمليات
على الأعمدة: فإضافة تركيبات مناسبة من العمودين $1, 2$ إلى العمودين
$3, 4$ تزيل $N$ \emph{عندما تكون $M$ قابلة للقلب} (بحلّ
$M\Lambda = -N$ من أجل معاملات التركيبة $\Lambda$)،
فيبقى $\det\begin{pmatrix} M & 0\\ 0 & P\end{pmatrix}$؛ ثم يعطي
\omterm{thm:b1:det:cofactor}{النشر بالعوامل المرافقة} على امتداد العمود الأول، مرتين، المقدارَ $\det
M\det P$ من أجل هذا الشكل الكتلي القطري. وإذا لم تكن $M$
قابلة للقلب، فأعمدتها مرتبطة، ومنه فالعمودان الأولان من
المصفوفة الكبيرة مرتبطان (لأن نصفيهما السفليين معدومان): فينعدم
الطرفان. وتمتدّ القاعدة إلى أيّ أحجام كتل بالحجة ذات
الحالتين نفسها --- وهي محرّك
\cref{exo:b1:det:10}.
\end{example}

\begin{example}[محدد $4 \times 4$، باستراتيجية]
\label{ex:b1:det:fourbyfour}
\[
\Delta = \begin{vmatrix}
1 & 2 & 3 & 4\\
2 & 3 & 4 & 1\\
3 & 4 & 1 & 2\\
4 & 1 & 2 & 3
\end{vmatrix}.
\]
مجموع كل سطر هو $10$: ومنه فالعملية $C_1 \leftarrow C_1 + C_2 +
C_3 + C_4$ تجعل العمود الأول ثابتًا، وبتعميل $10$
تبقى آحاد. ثم ينظّف $L_i \leftarrow L_i - L_1$ ($i \geq 2$)
العمودَ الأول:
\[
\Delta = 10\begin{vmatrix}
1 & 2 & 3 & 4\\
0 & 1 & 1 & -3\\
0 & 2 & -2 & -2\\
0 & -1 & -1 & -1
\end{vmatrix}
= 10\begin{vmatrix}
1 & 1 & -3\\
2 & -2 & -2\\
-1 & -1 & -1
\end{vmatrix}
= 10 \times 16 = 160,
\]
وينشر المحدد $3\times3$ الأخير على امتداد سطره الأول:
$1(2 - 2) - 1(-2 - 2) + (-3)(-2 - 2) = 0 + 4 + 12 = 16$. والعبرة:
عمليةٌ واحدة محسنة الاختيار (بملاحظة ثبات مجموع السطر) تغلب
ستة عشر عاملًا مرافقًا.
\end{example}

\begin{method}[اختيار استراتيجية للمحدد]
\label{met:b1:det:strategy}
امسح المصفوفة قبل حساب أيّ شيء.
\begin{enumerate}
  \item \emph{مجاميع سطور أو أعمدة ثابتة}: اجمع كل شيء في
        سطر واحد، وعمّل القيمة المشتركة
        (\cref{ex:b1:det:fourbyfour} و \cref{exo:b1:det:7}).
  \item \emph{بنية متكررة}: اطرح السطور أو الأعمدة
        المتجاورة لإنشاء أصفار؛ فتنهار الأنماط السلّمية
        نحو الشكل المثلثي، ومحدده يُقرأ على
        القطر.
  \item \emph{أصفار معزولة}: انشر على امتداد أخلى سطر
        (\cref{ex:b1:det:cofactor})؛ وتعطي العائلات
        التراجعية (ثلاثية الأقطار، \cref{exo:b1:det:6}) علاقات تراجعية
        بهذه الكيفية.
  \item \emph{وسيط}: يكون المحدد كثيرَ حدود
        فيه؛ فجد جذوره بملاحظة القيم المنحلّة
        (سطور متساوية، أعمدة متناسبة)، ثم ثبّت
        \omterm{def:b1:poly:def}{كثير الحدود} بالدرجة والمعامل المهيمن. ومن أجل
        مصفوفة \cref{exo:b1:det:7}: يعطي $m = 1$ ثلاثة
        سطور متساوية (بالرتبة $1$، وجذر مضاعف)، ويجعل $m = -2$
        مجموع السطور معدومًا (جذر آخر)؛ ودرجة المحدد
        $3$ في $m$ بحدّ مهيمن $-m^3$ (وهو
        جداء القطر المضاد $m\cdot m\cdot m$، وإشارته حسب
        ساروس $-1$)، ومنه يجب أن يكون $-(m+2)(m-1)^2$ --- دون أيّ
        نشر، وتتحقق الطريقتان إحداهما من الأخرى.
\end{enumerate}
\end{method}

\begin{example}[محدد فاندرموند]\label{ex:b1:det:vandermonde}
من أجل السلالم $x_1, \dots, x_n$:\index{محدد فاندرموند}
\[
V(x_1, \dots, x_n) =
\begin{vmatrix}
1 & x_1 & x_1^2 & \cdots & x_1^{n-1}\\
1 & x_2 & x_2^2 & \cdots & x_2^{n-1}\\
\vdots & & & & \vdots\\
1 & x_n & x_n^2 & \cdots & x_n^{n-1}
\end{vmatrix}
= \prod_{1 \leq i < j \leq n} (x_j - x_i) .
\]
وتخطيط البرهان (مفصَّلًا في \cref{exo:b1:det:5}): تنظّف عمليات
الأعمدة $C_k \leftarrow C_k - x_1 C_{k-1}$ من اليمين السطرَ
الأول، ويردّ تعميل كل سطر باقٍ الأمرَ إلى
$V(x_2, \dots, x_n)$. وهو غير معدوم إذا وفقط إذا كانت المقادير $x_i$ متمايزة مثنى مثنى
--- وهو المحدد وراء \omterm{thm:b1:poly:lagrange}{مقايسة لاغرانج}
(\cref{ex:b1:linmaps:interpolation}).
\end{example}

\section{الجمل الخطية}

\begin{definition}\label{def:b1:det:system}
الجملة الخطية\index{جملة خطية} ذات $n$ معادلة في $p$
مجهولًا هي $AX = B$ مع $A \in \mathcal{M}_{n,p}(K)$ و $B \in
K^n$؛ وتكون \emph{متجانسة} إذا كان $B = 0$. \omterm{def:b1:logic:sets}{ومجموعة} حلولها، إذا كانت
غير خالية، هي $X_0 + \ker A$: أي حلٌّ خاص زائد الحلّ المتجانس
العام --- وهي فضاء أفيني جزئي بُعده $p -
\operatorname{rk} A$ (بمبرهنة الرتبة).
\end{definition}

\begin{example}[البنية الأفينية، مجعولةً مرئية]
\label{ex:b1:det:affinestructure}
حُلَّ
\[
\begin{cases}
x + y + z = 3\\
x - y + 2z = 2 .
\end{cases}
\]
بطرح المعادلتين: $2y - z = 1$، ومنه $z = 2y - 1$ و $x =
3 - y - z = 4 - 3y$. وتكوّن الحلول المستقيمَ
\[
(x, y, z) = (4,\ 0,\ -1) + y\,(-3,\ 1,\ 2)
\qquad (y \in \R):
\]
أي الحلّ الخاص $X_0 = (4, 0, -1)$ (بالاختيار $y = 0$)
زائد مستقيم النواة $\ker A = \operatorname{Vect}(-3, 1, 2)$ للجملة
المتجانسة المرتبطة --- وللتحقق: $(-3) + 1 + 2 = 0$
و $(-3) - 1 + 4 = 0$. وهندسيًا، يتقاطع مستويان غير متوازيين من
$\R^3$ على امتداد مستقيم، وكان عدّ الأبعاد $p -
\operatorname{rk} A = 3 - 2 = 1$ يعلم ذلك قبل أن نحلّ
أيّ شيء. وتغيير الحلّ الخاص (ولنقل $y = 1$: $X_0' =
(1, 1, 1)$) يغيّر الوصف لا المستقيم: فللفضاء الأفيني
الجزئي مبادئ كثيرة واتجاه واحد.
\end{example}

\begin{theorem}[جمل كرامر المربعة]\label{thm:b1:det:cramer}
إذا كان $A \in GL_n(K)$، فللجملة $AX = B$ حلٌّ وحيد $X =
A^{-1}B$، وإحداثياته\index{قاعدة كرامر}
\[
x_j = \frac{\det A_j}{\det A},
\qquad A_j = A \text{ مع العمود } j \text{ مستبدَلًا بالمقدار } B .
\]
\end{theorem}

\begin{proof}
الوحدانية والوجود هما القابلية للقلب. وأمّا الصيغة:
فاكتب $B = \sum_k x_k C_k$ (بأعمدة $A$)؛ عندئذ، بتعدد
الخطية والتناوب،
\[
\det A_j = \det\Bigl(C_1, \dots, \sum_k x_k C_k, \dots, C_n\Bigr)
= \sum_k x_k \det(C_1, \dots, C_k, \dots, C_n)
= x_j \det A ,
\]
لأن لكل حدّ عدا $k = j$ عمودًا مكررًا.
\end{proof}

\begin{example}[كرامر بوسيط، كاملًا]
\label{ex:b1:det:cramerparameter}
من أجل $m \in \R$، حُلَّ
\[
\begin{cases}
x + m y = 1\\
m x + y = 2 .
\end{cases}
\]
المحدد هو $1 - m^2$. \emph{والحالة العامة} $m \neq \pm1$:
يعطي كرامر
\[
x = \frac{\begin{vmatrix} 1 & m\\ 2 & 1\end{vmatrix}}{1 - m^2}
= \frac{1 - 2m}{1 - m^2},
\qquad
y = \frac{\begin{vmatrix} 1 & 1\\ m & 2\end{vmatrix}}{1 - m^2}
= \frac{2 - m}{1 - m^2},
\]
أي حلًا نظيفًا واحدًا من أجل كل $m$ مقبول (وللتحقق عند $m = 0$:
$(1, 2)$، وهو صحيح بداهةً). \emph{والحالات المنحلّة}: عند $m = 1$
تُقرأ المعادلتان $x + y = 1$ و $x + y = 2$: فهما غير متوافقتين؛ وعند
$m = -1$ تُقرآن $x - y = 1$ و $-x + y = 2$، أي $x - y =
1$ و $x - y = -2$: فهما غير متوافقتين مرة أخرى. وانعدام
المحدد يعلن أن \emph{شيئًا} ما ينحلّ، لكنه
لا يقول أبدًا ماذا --- فالخالي أو غير المنتهي يجب أن يُفصل فيه بالنظر
في الطرف الأيمن. ولاحظ كذلك كيف تشير الصيغ إلى
حدودها هي: فعندما $m \to 1^{-}$، $x = \frac{1 - 2m}{1 - m^2} \to
-\infty$؛ وتهرب نقطة الحلّ كلما صار المستقيمان
متوازيين.
\end{example}

\begin{method}[إزاحة غاوس على الجمل]\label{met:b1:det:gauss}
قلّص بالسطور المصفوفةَ الموسَّعة $(A \mid B)$ إلى الصورة
الدرجية.\index{إزاحة غاوس}
\begin{enumerate}
  \item إذا ظهر محور في العمود الأخير (سطر $0 = 1$): فلا
        حلّ.
  \item وإلا فتنقسم المجاهيل إلى \emph{مجاهيل محورية} و
        \emph{مجاهيل \omterm{def:b1:vspaces:free}{حرة}} (وسائط)؛ ويعبّر التعويض الرجعي
        عن الأولى بالثانية: فتكون \omterm{def:b1:logic:sets}{مجموعة} الحلول فضاءً
        أفينيًا جزئيًا بُعده $=$ عدد المجاهيل \omterm{def:b1:vspaces:free}{الحرة}.
\end{enumerate}
وصيغ كرامر للنظرية وللجمل الصغيرة؛ والإزاحة هي
الخوارزمية العملية.
\end{method}

\begin{example}[مناقشة بوسيط]\label{ex:b1:det:parameter}
من أجل $m \in \R$، انظر في
\[
\begin{cases}
x + y + mz = 1\\
x + my + z = 1\\
mx + y + z = 1 .
\end{cases}
\]
للمصفوفة محدد $-(m+2)(m-1)^2$ (محسوبًا في
\cref{exo:b1:det:7} بإضافة كل الأعمدة إلى الأول). ومن أجل $m
\neq 1, -2$:
حلٌّ وحيد $x = y = z = \frac{1}{m+2}$ (بالتناظر). ومن أجل $m =
1$: معادلة واحدة مكررة ثلاث مرات، أي مستوي حلول. ومن أجل $m =
-2$: يعطي جمع المعادلات الثلاث $0 = 3$، فلا حلّ.
\end{example}

\begin{remark}[مزالق شائعة]
\label{rem:b1:det:pitfalls}
\emph{المحدد ليس \omterm{def:b1:linmaps:def}{خطيًا} في المصفوفة}: $\det(A + B)
\neq \det A + \det B$ (بل أصلًا $\det(I_2 + I_2) = 4 \neq 2$)؛ فهو
خطي في كل \emph{عمود} على حدة، وهذا شيء مختلف
تمامًا.
\emph{السلّمة}: $\det(\lambda A) = \lambda^n\det A$، لا
$\lambda\det A$ --- لأن كل عمود من الأعمدة $n$ يُضرب في سلّم.
\emph{ليست كل \omterm{met:b1:matrices:gauss}{العمليات على السطور} مجانية}: فالعملية $L_i \leftarrow L_i +
\lambda L_j$ تحفظ المحدد، لكن المبادلة تغيّر
الإشارة و $L_i \leftarrow \lambda L_i$ يضربه في $\lambda$
--- وأخطاء مسك الدفاتر هنا هي المصدر الكلاسيكي للإشارات
الخاطئة في الحسابات القائمة على الإزاحة.
\emph{المحدد المعدوم بداية لا نهاية}: فهو يقول
«الرتبة $< n$» لكن لا يقول أيّ رتبة؛ ولا يحدّدها إلا عمل إضافي (الصورة
الدرجية، أو المحددات الجزئية في \cref{exo:b1:det:12}) --- قارن
حالة $m = 1$ إزاء $m = -2$ في \cref{ex:b1:det:parameter}.
\emph{كرامر يحتاج إلى القابلية للقلب}: فعندما $\det A = 0$ تكون
الصيغ $x_j = \det A_j/\det A$ بلا معنى، وقد يكون للجملة
حلول (لا تُحصى) تمامًا.
\emph{المصفوفات المربعة وحدها لها محددات}: فمن أجل جملة
مستطيلة، تكون الإزاحة الأداةَ الوحيدة.
\end{remark}

\begin{remark}[إلى أين تذهب المحددات]
\label{rem:b1:det:whereused}
ثلاث حيوات تنتظر هذا السلّم. \emph{هندسية}: فالمقدار $\abs{\det}$ هو
عامل تسليم المساحة أو الحجم \omterm{def:b1:logic:map}{للتطبيق} المرتبط --- وهو
مدقَّق من أجل المستوي في \cref{ch:b1:euclid} و، بوصفه يعقوبيَّ
تغيير متغيرات، في التكاملات المضاعفة في مجلّد السنة~الثانية.
\emph{وجبرية}: فالمقدار $\det(A - \lambda I)$، وهو \omterm{def:b1:poly:def}{كثير الحدود}
المميّز، يفتح نظرية القيم الذاتية في السنة~الثانية ---
والمتطابقة $A^2 - (\operatorname{tr} A)A + (\det A)I = 0$ في
مسألة نهاية الأسبوع في \cref{ch:b1:matrices} ظلُّه الأول.
\emph{وتحليلية}: فمحددات مصفوفات خاصة (فاندرموند و
كوشي وغرام) تفصل متى تكون مسائل المقايسة والتفكيك و
\omterm{def:b1:linmaps:projection}{الإسقاط} مطروحة طرحًا حسنًا؛ وتقوّم مسألة نهاية الأسبوع أدناه
العائلتين الأوليين تقويمًا كاملًا.
\end{remark}

\begin{remark}[منظورات داخل الكتاب 3]
\label{rem:b1:det:perspectives}
يغلق هذا الفصل عمود الجبر الخطي في المجلّد، و
يصرف فصلاه الباقيان الأرباح. ففي
\cref{ch:b1:euclid}: تختبر مصفوفة غرام
$\bigl(\langle v_i, v_j\rangle\bigr)$ الحريةَ بمحدد
(\cref{exo:b1:euclid:11})، وتنقسم تقايسات المستوي
إلى دورانات وانعكاسات حسب إشارة
محددها --- ويجري تصنيف مسألة نهاية الأسبوع هناك عليها. وفي
\cref{ch:b1:multivar}: يكون مقدار مونج $rt -
s^2$ محددَ المصفوفة المتماثلة للمشتقات
الثانية، وتكون المعادلات الناظمية للمربعات الصغرى جملةَ
كرامر مصفوفتها مصفوفة غرام (ومنه مصفوفة عزوم) ---
قابلةً للقلب بالضبط بالمحكات ذات النكهة الفاندرموندية
المؤسَّسة هنا. وحين تؤكد تلك الفصول «قابلة للقلب» أو
«موجبة»، فالإيصالات في هذا الفصل.
\end{remark}

\section{تمارين}

\begin{exercise}[$\star$]\label{exo:b1:det:1}
احسب:
\[
\begin{vmatrix} 3 & 1\\ 5 & 2 \end{vmatrix},
\qquad
\begin{vmatrix} 1 & 2 & 3\\ 4 & 5 & 6\\ 7 & 8 & 9\end{vmatrix},
\qquad
\begin{vmatrix} 1 & 1 & 1\\ 1 & 2 & 4\\ 1 & 3 & 9\end{vmatrix}.
\]
\end{exercise}

\begin{exercise}[$\star$]\label{exo:b1:det:2}
من أجل أيّ $\lambda \in \R$ تكون العائلة $\bigl((1, 1, \lambda),
(1, \lambda, 1), (\lambda, 1, 1)\bigr)$ \omterm{def:b1:vspaces:free}{أساسًا} للمقدار $\R^3$؟
\end{exercise}

\begin{exercise}[$\star$]\label{exo:b1:det:3}
حُلَّ بقاعدة كرامر:
\[
\begin{cases}
2x + y = 5\\
3x - 2y = 4 ,
\end{cases}
\qquad\text{ثم}\qquad
\begin{cases}
x + y + z = 6\\
x - y + z = 2\\
2x + y - z = 1 .
\end{cases}
\]
\end{exercise}

\begin{exercise}[$\star$]\label{exo:b1:det:4}
حُلَّ \omterm{met:b1:det:gauss}{بإزاحة غاوس}، واصفًا \omterm{def:b1:logic:sets}{مجموعة} الحلول:
\[
\begin{cases}
x + 2y - z + t = 1\\
2x + 4y + z - t = 5\\
x + 2y + 2z - 2t = 4 .
\end{cases}
\]
\end{exercise}

\begin{exercise}[$\star\star$]\label{exo:b1:det:5}
برهن على صيغة فاندرموند في \cref{ex:b1:det:vandermonde}
بالاستقراء على $n$، مع عمليات الأعمدة $C_k \leftarrow C_k -
x_1 C_{k-1}$ منفَّذةً من $k = n$ نزولًا إلى $k = 2$.
\end{exercise}

\begin{exercise}[$\star\star$]\label{exo:b1:det:6}
(ثلاثية الأقطار) ليكن $D_n$ المحدد $n \times n$ ذا $2$ على
القطر و $1$ على القطرين المجاورين و $0$ فيما عداهما.
وبالنشر على امتداد السطر الأول، برهن على $D_n = 2D_{n-1} - D_{n-2}$ و
احسب $D_n$ ($D_1 = 2$ و $D_2 = 3$).
\end{exercise}

\begin{exercise}[$\star\star$]\label{exo:b1:det:7}
أكمل \cref{ex:b1:det:parameter}: احسب المحدد
$\begin{vmatrix} 1 & 1 & m\\ 1 & m & 1\\ m & 1 & 1\end{vmatrix}$
بالعملية $C_1 \leftarrow C_1 + C_2 + C_3$، ونفّذ
المناقشة الكاملة للجملة.
\end{exercise}

\begin{exercise}[$\star\star$]\label{exo:b1:det:8}
لتكن $A \in \mathcal{M}_n(\R)$ ذات مركّبات \emph{صحيحة}. برهن
على أن للمقدار $A$ معكوسًا بمركّبات صحيحة إذا وفقط إذا كان $\det A
= \pm 1$. \emph{(ومن أجل الاتجاه المباشر، خذ المحددات؛ ومن أجل
العكس، اقبل --- أو برهن من أجل $n \leq 3$ عبر العوامل المرافقة --- أن
$A^{-1} = \frac{1}{\det A}\,\operatorname{Com}(A)^{\mathsf T}$ بمصفوفة عوامل
مرافقة صحيحة.)}
\end{exercise}

\begin{exercise}[$\star\star\star$]\label{exo:b1:det:9}
احسب المحدد $n \times n$ للمصفوفة $aI + bJ$
(\cref{exo:b1:matrices:9})، أي ذات $a + b$ على القطر و
$b$ فيما عداه. \emph{(أضف كل الأعمدة إلى الأول، وعمّل، ثم
نظّف.)} واستعد شرط القابلية للقلب $a \neq 0$ و $a + nb
\neq 0$.
\end{exercise}

\begin{exercise}[$\star\star\star$]\label{exo:b1:det:10}
لتكن $A, B \in \mathcal{M}_n(\R)$. برهن على أن
\[
\det\begin{pmatrix} A & B\\ B & A \end{pmatrix}
= \det(A + B)\,\det(A - B),
\]
بعمليات كتلية على الأعمدة والسطور ($C_1 \leftarrow C_1 + C_2$، ثم
$L_2 \leftarrow L_2 - L_1$، في صورة كتلية)، مفترضًا قاعدة الكتل
المثلثية الطبيعية $\det\begin{pmatrix} M & N\\ 0 &
P\end{pmatrix} = \det M \det P$ --- المبرهن عليها من أجل كتل $2 \times 2$
في \cref{ex:b1:det:blocktriangular}.
\end{exercise}

\begin{exercise}[$\star\star$]\label{exo:b1:det:11}
(الدائرية من الرتبة $3$) ليكن $a, b, c \in \C$ و
\[
\Delta = \begin{vmatrix}
a & b & c\\
c & a & b\\
b & c & a
\end{vmatrix}.
\]
برهن على أن $\Delta = (a + b + c)(a^2 + b^2 + c^2 - ab - bc -
ca)$، وعمّل تعميلًا كاملًا على $\C$ باستعمال $j =
\eu^{2\iu\pi/3}$:
\[
\Delta = (a + b + c)(a + jb + j^2c)(a + j^2b + jc) .
\]
\emph{(ابدأ بالمقدار $C_1 \leftarrow C_1 + C_2 + C_3$؛ ومن أجل
الصورة العقدية، لاحظ أن العمود $(1, j, j^2)^{\mathsf T}$
يكاد يسلك سلوك متجهة ذاتية.)}
\end{exercise}

\begin{exercise}[$\star\star\star$]\label{exo:b1:det:12}
(الرتبة والمحددات الجزئية) لتكن $A \in \mathcal{M}_{n,p}(K)$. برهن على أن
$\operatorname{rk} A$ يساوي أكبر حجم $r$ لمصفوفة جزئية
$r \times r$ قابلة للقلب من $A$ (والمصفوفة الجزئية تحتفظ
بالمركّبات عند تقاطعات $r$ سطرًا مختارًا و $r$ عمودًا
مختارًا). \emph{(إذا كان $\operatorname{rk} A = r$، فاختر $r$ عمودًا
حرًا، ثم $r$ سطرًا حرًا من الكتلة $n \times r$ الناتجة؛
وبالعكس، تفرض مصفوفة جزئية قابلة للقلب أن تكون الأعمدة
المقابلة من $A$ \omterm{def:b1:vspaces:free}{حرة}.)}
\end{exercise}

\section{مسألة: المتناوب المزدوج لكوشي}

\begin{problem}\label{pb:b1:det:1}
محددان يحكمان تطبيقات هذا الفصل: \omterm{ex:b1:det:vandermonde}{محدد
فاندرموند}، المقوَّم في \cref{exo:b1:det:5}، و
\emph{\omterm{pb:b1:det:1}{محدد كوشي}} $\det\bigl(\frac{1}{a_i +
b_j}\bigr)$، المقوَّم هنا. وحول هذين يجمع هذا العمل
عدّةَ المتناوبات: حيل الأعمدة الكثيرةالحدود، والمقايسة
بكرامر، \omterm{pb:b1:det:1}{ومصفوفة هيلبرت}، ومميّز \omterm{def:b1:poly:def}{كثير حدود} تكعيبي، و
طريقة كثيرات الحدود المتناوبة. وفي كل ما يأتي، يرمز $V(x_1,
\dots, x_n) = \prod_{i < j}(x_j - x_i)$ إلى قيمة
فاندرموند.\index{محدد كوشي}\index{مصفوفة هيلبرت}
\index{مميّز}

\medskip
\noindent\textbf{الجزء 1 --- عدّة فاندرموند.}
\begin{enumerate}
  \item احسب $V(1, 2, 3, 4)$، وتذكّر لماذا تكون المقايسة عند
        $n$ عقدة متمايزة مثنى مثنى جملةَ كرامر.
  \item (المتناوب الكثيرالحدود) ليكن $P_0, \dots, P_{n-1}$
        \emph{واحديات} مع $\deg P_k = k$. برهن على
        \[
        \det\bigl(P_{j-1}(x_i)\bigr)_{1 \leq i, j \leq n}
        = V(x_1, \dots, x_n) :
        \]
        فعمليات الأعمدة تستبدل بكل عمود قوى أيَّ سلّم
        واحديّ، بالمجّان.
  \item طبّق السؤال 2 على كثيرات الحدود الثنائية $B_k =
        \frac{X(X-1)\cdots(X-k+1)}{k!}$: وبرهن على أنه من أجل
        \emph{أعداد صحيحة} $m_1 < m_2 < \dots < m_n$،
        \[
        \frac{V(m_1, \dots, m_n)}{0!\,1!\,2!\cdots(n-1)!}
        \in \N :
        \]
        أي إن جداء كل الفروق مثنى مثنى لعدد $n$ من الأعداد الصحيحة
        يقبل القسمة على العاملي الفائق $0!\,1!\cdots(n-1)!$.
  \item برهن على $\det\bigl(x_i^{\,j}\bigr)_{1 \leq i, j \leq n} =
        x_1 x_2 \cdots x_n\, V(x_1, \dots, x_n)$ (والقوى تبدأ
        الآن من $1$).
  \item (مصفوفة العزوم) لتكن $S = \bigl(p_{i+j-2}\bigr)_{1 \leq
        i, j \leq n}$ حيث $p_k = x_1^k + \dots + x_n^k$. برهن
        على أن $S = W^{\mathsf T} W$ من أجل المصفوفة $W =
        (x_i^{\,j-1})_{ij}$، واستنتج
        \[
        \det S = V(x_1, \dots, x_n)^2 ,
        \]
        واختم بأن: $n$ عددًا \emph{حقيقيًا} تكون متمايزة مثنى
        مثنى إذا وفقط إذا كانت مصفوفة عزومها
        قابلة للقلب، وأن $\det S \geq 0$ دائمًا.
\end{enumerate}

\medskip
\noindent\textbf{الجزء 2 --- المقايسة، مُعادًا النظر فيها.} العقد
$x_1 < \dots < x_n$، والقيم $y_1, \dots, y_n$.
\begin{enumerate}[resume]
  \item اكتب الشروط «$P = c_0 + c_1X + \dots +
        c_{n-1}X^{n-1}$ يقايس» جملةً خطية في
        المقادير $c_k$ مصفوفتها $W$، واستعد من $\det W = V \neq
        0$ وجود المقايِس ووحدانيته
        (وقارن البرهانين السابقين،
        \cref{thm:b1:poly:lagrange} و
        \cref{ex:b1:linmaps:interpolation}).
  \item بقاعدة كرامر وبالنشر بالعوامل المرافقة للمحدد المعني
        على امتداد عموده الأخير، برهن على أن
        المعامل المهيمن للمقايِس هو
        \[
        c_{n-1} = \sum_{i=1}^{n}
        \frac{y_i}{\prod_{j \neq i}(x_i - x_j)} .
        \]
  \item (فاندرموند الملتحم) احسب
        \[
        \begin{vmatrix}
        1 & x_1 & x_1^2\\
        0 & 1 & 2x_1\\
        1 & x_2 & x_2^2
        \end{vmatrix}
        = (x_2 - x_1)^2 ,
        \]
        وفسّر: فالمعطيات $\bigl(P(x_1), P'(x_1),
        P(x_2)\bigr)$ تحدّد $P \in \R_2[X]$ وحيدًا عندما يكون
        $x_1 \neq x_2$ (وهي مقايسة إرميت).
  \item جد $P \in \R_2[X]$ الوحيد الذي يحقق $P(0) = 1$ و $P'(0)
        = 0$ و $P(1) = 2$، وتحقق من جوابك إزاء
        السؤال 8.
\end{enumerate}

\medskip
\noindent\textbf{الجزء 3 --- \omterm{pb:b1:det:1}{محدد كوشي}.} ليكن $a_1,
\dots, a_n$ و $b_1, \dots, b_n$ سلالم تحقق $a_i + b_j
\neq 0$ من أجل كل $i, j$، ولتكن
\[
C_n = \det\Bigl(\frac{1}{a_i + b_j}\Bigr)_{1 \leq i, j \leq n} .
\]
\begin{enumerate}[resume]
  \item احسب $C_2$ باليد وضعه في الصورة
        «جداءات فروق على جداءات مجاميع».
  \item من أجل $n \geq 2$، نفّذ $L_i \leftarrow L_i - L_n$ ($i <
        n$) وعمّل السطور والأعمدة لتبرهن على
        \[
        C_n = \frac{\prod_{i<n}(a_n - a_i)}{\prod_{j}(a_n +
        b_j)}\;\det M,
        \]
        حيث تتفق $M$ مع مصفوفة كوشي على السطور $i < n$
        ولها السطر الأخير $(1, 1, \dots, 1)$.
  \item نفّذ $C_j \leftarrow C_j - C_n$ ($j < n$) على $M$، و
        عمّل مرة أخرى، واختم بالاستقراء
        \emph{المتناوب المزدوج لكوشي}:
        \[
        C_n = \frac{\prod_{1 \leq i < j \leq n}(a_j - a_i)(b_j -
        b_i)}{\prod_{i, j}(a_i + b_j)} .
        \]
  \item استنتج محك القابلية للقلب (المقادير $a_i$ متمايزة مثنى
        مثنى والمقادير $b_j$ متمايزة مثنى مثنى). ومن أجل
        \emph{\omterm{pb:b1:det:1}{مصفوفة هيلبرت}} $H_n = \bigl(\frac{1}{i + j -
        1}\bigr)$: احسب $\det H_2$ و $\det H_3$ من
        الصيغة، وتحقق من أن $H_2^{-1}$ ذات مركّبات صحيحة.
  \item بيّن أنه من أجل $b_j$ متمايزة مثنى مثنى وأيّ طرف
        أيمن، يكون للجملة $\sum_j \frac{c_j}{a_i + b_j} = y_i$
        ($i = 1, \dots, n$) حلٌّ وحيد، واربط
        ذلك بوجود التفكيك إلى عناصر بسيطة ذات \omterm{def:b1:fractions:field}{أقطاب} بسيطة
        ووحدانيته
        (\cref{thm:b1:fractions:complex}).
\end{enumerate}

\medskip
\noindent\textbf{الجزء 4 --- مميّز \omterm{def:b1:poly:def}{كثير حدود} تكعيبي.} ليكن
$\lambda_1, \lambda_2, \lambda_3$ جذور $X^3
+ pX + q$ (في $\C$)، ولتكن $p_k = \lambda_1^k + \lambda_2^k + \lambda_3^k$.
\begin{enumerate}[resume]
  \item باستعمال $\lambda^3 = -p\lambda - q$ عند كل جذر و
        فييت ($p_1 = 0$)، احسب $p_2 = -2p$ و $p_3 = -3q$ و
        $p_4 = 2p^2$.
  \item مع السؤال 5 (على $\C$، مع الاحتفاظ بالمقدار $\det S = V^2$)،
        احسب
        \[
        \operatorname{disc} = V(\lambda_1, \lambda_2,
        \lambda_3)^2 = \begin{vmatrix}
        3 & 0 & -2p\\
        0 & -2p & -3q\\
        -2p & -3q & 2p^2
        \end{vmatrix}
        = -4p^3 - 27q^2 .
        \]
  \item استنتج: \omterm{def:b1:poly:def}{لكثير الحدود} $X^3 + pX + q$ جذر مضاعف إذا وفقط
        إذا كان $4p^3 + 27q^2 = 0$؛ وتحقق على $X^3 - 3X + 2 = (X -
        1)^2(X + 2)$.
  \item افترض $p, q$ حقيقيين. برهن على أن للتكعيبي ثلاثة
        جذور حقيقية متمايزة إذا وفقط إذا كان $\operatorname{disc}
        > 0$، وجذرًا حقيقيًا واحدًا مع جذرين عقديين مرافقين إذا
        وفقط إذا كان $\operatorname{disc} < 0$. \emph{(إذا كان
        $\lambda_3 = \conj{\lambda_2} \neq \lambda_2$ و
        $\lambda_1 \in \R$، فبيّن أن $V$ تخيّلي محض.)}
\end{enumerate}

\medskip
\noindent\textbf{الجزء 5 --- الأرباح، وطريقة
التناوب.}
\begin{enumerate}[resume]
  \item من أجل $0 < a_1 < a_2 < \dots < a_n$، بيّن
        $\det\bigl(\frac{1}{a_i + a_j}\bigr) > 0$.
  \item احسب $\det\bigl(\binom{m_i}{j-1}\bigr)_{1 \leq i, j
        \leq 3}$ من أجل $(m_1, m_2, m_3) = (2, 4, 7)$، أولًا
        بالسؤالين 2--3، ثم بالنشر المباشر.
  \item ليكن $\lambda_1, \dots, \lambda_n$ متمايزة مثنى مثنى
        وغير معدومة. وباستعمال مصفوفة فاندرموند قابلة للقلب،
        برهن مرة أخرى على أن المتتاليات الهندسية
        $\bigl((\lambda_i^{\,k})_{k \geq 0}\bigr)_{1 \leq i \leq
        n}$ تكوّن \omterm{def:b1:vspaces:free}{عائلة حرة} من فضاء المتتاليات.
  \item احسب $\det\bigl(\frac{1}{i + j}\bigr)_{1 \leq i, j
        \leq 3}$ من المتناوب المزدوج.
  \item (كثيرات الحدود المتناوبة) سمِّ \omterm{def:b1:poly:def}{كثير حدود} $F$ في $x_1,
        \dots, x_n$ \emph{متناوبًا} إذا كانت مبادلة أيّ متغيرين
        تغيّر إشارته. بيّن أن \omterm{def:b1:poly:def}{كثير الحدود} المتناوب $F$
        ينعدم كلما كان $x_i = x_j$ ($i \neq j$)، واستنتج
        --- متغيرًا في كل مرة، بمبرهنة العامل ---
        أن $F$ يقبل القسمة على $\prod_{i<j}(x_j - x_i)$.
  \item استعمل السؤال 23 لتعيد البرهان على صيغة فاندرموند
        دون أيّ استقراء: فالمحدد $\det(x_i^{\,j-1})$ كثيرُ
        حدود متناوب درجته الكلية $\binom n2$،
        ومنه فهو مضاعفٌ \emph{ثابت} للمقدار $\prod_{i<j}(x_j -
        x_i)$؛ وعيّن الثابت بمقارنة وحيد حدّ واحد.
  \item توليفة، في أربع جمل: أيّ خاصية وحيدة من خصائص
        المحدد (أيّ بديهية) تولّد كل
        التعميلات في هذه المسألة؛ ولماذا تحوّل متطابقة مصفوفة
        العزوم في السؤال 5 عبارةً عن التمايز
        \emph{العقدي} إلى اختبار إشارة \emph{حقيقي}
        قابل للحساب؛ وأيّ مصفوفتين كلاسيكيتين قُوّمتا
        تقويمًا كاملًا هنا وأيّ مسائل خطية تحكمانها؛ وكيف
        تفسّر طريقة التناوب في السؤالين
        23--24، دفعة واحدة، سببَ استمرار ظهور $\prod_{i<j}(x_j -
        x_i)$. وسمِّ مبرهنة الجزء 3.
\end{enumerate}
\end{problem}
